% =============================================================================
% IEEEtran Journal Template
% Engineering Paper Writing Workflow v1.1.0
%
% TARGET JOURNALS — select one in \documentclass options:
%   TIE   : IEEE Transactions on Industrial Electronics
%   TIM   : IEEE Transactions on Instrumentation and Measurement
%   TMECH : IEEE/ASME Transactions on Mechatronics
%   TAC   : IEEE Transactions on Automatic Control
%   TRO   : IEEE Transactions on Robotics
%   RAL   : IEEE Robotics and Automation Letters
%
% JOURNAL-SPECIFIC NOTES:
%   TIE   : Max 8 pages (regular), color figures OK in online version
%   TIM   : Max 10 pages, measurement focus required in abstract
%   TMECH : Max 10 pages (regular), must include system implementation
%   TAC   : Max 8 pages (regular), 4 pages (technical notes)
%           Use \documentclass[technote]{IEEEtran} for technical notes
%   TRO   : Max 12 pages, video attachment recommended
%   RAL   : Max 8 pages, presented at ICRA or IROS, strict deadline
%
% COMPILE SEQUENCE:
%   pdflatex ieee_template.tex
%   bibtex   ieee_template
%   pdflatex ieee_template.tex
%   pdflatex ieee_template.tex
% =============================================================================

% ---------------- Document Class --------------------------------------------
% Options: journal (default) | technote | peerreview | peerreviewca
% Two-column is the default for journal; comment out 'twocolumn' if needed
\documentclass[journal]{IEEEtran}
% For TAC technical notes:  \documentclass[technote]{IEEEtran}
% For peer-review version:  \documentclass[journal,peerreview]{IEEEtran}

% ---------------- Packages --------------------------------------------------
% Mathematics
\usepackage{amsmath}          % align, equation, cases, etc.
\usepackage{amssymb}          % \mathbb, \mathcal, etc.
% NOTE: amsthm is intentionally NOT loaded.
% IEEEtran defines its own 'proof' environment; loading amsthm causes
% a "proof already defined" conflict. Use \newtheorem below for theorem-type
% environments, and IEEEtran's built-in proof environment directly.
% \usepackage{amsthm}  ← DO NOT uncomment
\usepackage{bm}               % \bm{} for bold math (preferred over \mathbf)

% Graphics
\usepackage{graphicx}         % \includegraphics
\usepackage{epstopdf}         % convert .eps to .pdf automatically
% \usepackage{subfig}         % sub-figures (use if needed)

% Tables
\usepackage{booktabs}         % \toprule, \midrule, \bottomrule
\usepackage{multirow}         % \multirow{}

% Algorithms
\usepackage{algorithm}
\usepackage{algpseudocode}    % \Procedure, \State, \For, \If, etc.

% Units (SI)
\usepackage{siunitx}          % \SI{9.81}{\meter\per\second\squared}

% Hyperlinks (comment out for final camera-ready if journal disallows)
\usepackage[colorlinks=true, linkcolor=blue, citecolor=blue, urlcolor=blue]{hyperref}

% Cross-referencing
\usepackage{cleveref}         % \cref{fig:block} → "Fig. 1"

% Misc
\usepackage{cite}             % compact citation ranges [1]–[3]
\usepackage{url}              % \url{}
\usepackage{balance}          % balance last page columns (add \balance before \bibliography)

% ---------------- Theorem Environments --------------------------------------
\newtheorem{theorem}{Theorem}
\newtheorem{lemma}[theorem]{Lemma}
\newtheorem{corollary}[theorem]{Corollary}
\newtheorem{proposition}[theorem]{Proposition}
\newtheorem{remark}{Remark}
\newtheorem{assumption}{Assumption}
\newtheorem{definition}{Definition}

% ---------------- Custom Commands -------------------------------------------
% Notation shortcuts — keep consistent with method/notation.md
\newcommand{\bx}{\bm{x}}          % state vector
\newcommand{\bu}{\bm{u}}          % control input
\newcommand{\by}{\bm{y}}          % output
\newcommand{\bA}{\bm{A}}          % system matrix
\newcommand{\bB}{\bm{B}}          % input matrix
\newcommand{\bC}{\bm{C}}          % output matrix
\newcommand{\bK}{\bm{K}}          % gain matrix
\newcommand{\bP}{\bm{P}}          % Lyapunov / Riccati matrix
\newcommand{\bQ}{\bm{Q}}          % weight matrix Q
\newcommand{\bR}{\bm{R}}          % weight matrix R
\newcommand{\R}{\mathbb{R}}        % real numbers
\newcommand{\norm}[1]{\left\|#1\right\|}
\newcommand{\abs}[1]{\left|#1\right|}
\newcommand{\T}{^{\top}}           % transpose

% =============================================================================
% DOCUMENT START
% =============================================================================
\begin{document}

% ---------------- Title ------------------------------------------------------
\title{%
  % TODO: Replace with your paper title
  % IEEE style: avoid abbreviations in title, no math symbols
  Insert Your Paper Title Here%
  % Example: "Robust Preview Control for a Six-Degree-of-Freedom
  %           Motion Simulation Platform via Horizon Optimization"
}

% ---------------- Authors ----------------------------------------------------
% IEEE format: \IEEEauthorblockN for names, \IEEEauthorblockA for affiliation
\author{%
  \IEEEauthorblockN{Hungsun Son,~\IEEEmembership{Member,~IEEE,}
                    and Author B}
  \IEEEauthorblockA{%
    Department of Mechanical Engineering,\\
    Ulsan National Institute of Science and Technology (UNIST),\\
    Ulsan 44919, Republic of Korea\\
    Email: hson@unist.ac.kr%
  }%
  % Add more authors:
  % \and
  % \IEEEauthorblockN{Author C}
  % \IEEEauthorblockA{Affiliation C \\ email@c.edu}
}

% Marks are for: Member, Senior Member, Fellow, Life Member, Life Fellow
% If not a member: remove ,~\IEEEmembership{...}

% ---------------- Running Header (for peer-review) --------------------------
% \markboth{IEEE TRANSACTIONS ON ..., VOL. XX, NO. X, MONTH YEAR}%
% {Son \MakeLowercase{\textit{et al.}}: Short Paper Title}

\maketitle

% ---------------- Abstract --------------------------------------------------
\begin{abstract}
% IEEE abstract guidelines:
%   - 150–250 words
%   - No citations, no math notation if avoidable, no abbreviations on first use
%   - State: (1) problem/motivation, (2) proposed method, (3) key results, (4) significance
%
% TIM note: First sentence must mention the measurement/instrumentation aspect.
%
% Structure template:
%   [Problem]  This paper addresses the problem of ...
%   [Gap]      Existing approaches suffer from ... because ...
%   [Method]   We propose a ... that ...
%   [Result]   Experiments/simulations demonstrate ... improvement of X% over ...
%   [Impact]   The proposed method enables ...
%
TODO: Write abstract here. (150--250 words)
\end{abstract}

% ---------------- Index Terms -----------------------------------------------
\begin{IEEEkeywords}
% 4–8 keywords; use IEEE Taxonomy terms where possible
% https://www.ieee.org/publications/services/thesaurus.html
keyword1, keyword2, keyword3, keyword4
\end{IEEEkeywords}

% =============================================================================
% SECTION I — INTRODUCTION
% =============================================================================
\section{Introduction}
\label{sec:intro}
% IEEE Introduction structure:
%   Paragraph 1: Problem statement and motivation (why does this matter?)
%   Paragraph 2–3: Survey of existing work (what has been done?)
%   Paragraph 4: Gap / limitation of existing work (what is still missing?)
%   Paragraph 5: Proposed approach and novelty (what do WE do?)
%   Paragraph 6: Paper organization (brief, e.g., "Section II presents...")
%
% Novelty claims go here as explicit statements:
%   "The main contributions of this paper are as follows:"
%   1) ...
%   2) ...
%   3) ...

TODO: Write introduction.

The main contributions of this paper are summarized as follows:
\begin{enumerate}
  \item TODO: Contribution 1 (linked to simulation/comparison result)
  \item TODO: Contribution 2 (linked to experiment result)
  \item TODO: Contribution 3 (if applicable)
\end{enumerate}

The remainder of this paper is organized as follows.
Section~\ref{sec:related} reviews related work.
Section~\ref{sec:problem} formulates the problem.
Section~\ref{sec:method} presents the proposed method.
Section~\ref{sec:simulation} validates the method via simulation.
Section~\ref{sec:experiment} presents experimental results.
Section~\ref{sec:conclusion} concludes the paper.


% =============================================================================
% SECTION II — RELATED WORK (or BACKGROUND / PRELIMINARIES)
% =============================================================================
\section{Related Work}
\label{sec:related}
% Organize into 2–4 thematic paragraphs. Each paragraph covers one research thread.
% End with a paragraph that explicitly states the gap that your paper fills.
% Every statement must be supported by at least one citation [X].

TODO: Write related work.


% =============================================================================
% SECTION III — PROBLEM FORMULATION
% =============================================================================
\section{Problem Formulation}
\label{sec:problem}

\subsection{System Model}
% Define your system with numbered equations.
% ALL notation must appear in method/notation.md before being used here.

Consider the discrete-time LTI system:
\begin{equation}
  \bx_{k+1} = \bA\,\bx_k + \bB\,\bu_k,
  \label{eq:system}
\end{equation}
\begin{equation}
  \by_k = \bC\,\bx_k,
  \label{eq:output}
\end{equation}
where $\bx_k \in \R^n$ is the state vector, $\bu_k \in \R^m$ is the control input,
$\by_k \in \R^p$ is the measured output, and
$\bA \in \R^{n \times n}$, $\bB \in \R^{n \times m}$, $\bC \in \R^{p \times n}$
are the system, input, and output matrices, respectively.

\begin{assumption}
  \label{ass:stabilizable}
  The pair $(\bA, \bB)$ is stabilizable and $(\bA, \bC)$ is detectable.
\end{assumption}

\subsection{Control Objective}
% State the control objective as a formal optimization or requirement.

TODO: State the control objective.

The control objective is to find a control law $\bu_k = \pi(\bx_k)$ such that:
\begin{equation}
  \min_{\bu_k} \sum_{k=0}^{\infty}
    \left( \bx_k\T \bQ\,\bx_k + \bu_k\T \bR\,\bu_k \right),
  \label{eq:objective}
\end{equation}
subject to \eqref{eq:system}--\eqref{eq:output} and Assumption~\ref{ass:stabilizable}.


% =============================================================================
% SECTION IV — PROPOSED METHOD
% =============================================================================
\section{Proposed Method}
\label{sec:method}

\subsection{Overview}
% One paragraph overview with a block diagram (Fig.~\ref{fig:block}).

TODO: Describe method overview.

\cref{fig:block} illustrates the overall architecture of the proposed system.

\begin{figure}[htbp]
  \centering
  % Replace with your actual figure file (PDF or EPS preferred for IEEEtran)
  % \includegraphics[width=\columnwidth]{figures/block_diagram.pdf}
  \rule{\columnwidth}{0.4pt}   % placeholder — remove when figure is ready
  \vspace{1em}
  \caption{Block diagram of the proposed control system.
           TODO: Add meaningful caption.}
  \label{fig:block}
\end{figure}

\subsection{Algorithm Design}
% Present the core algorithm.

TODO: Describe the algorithm.

\begin{algorithm}[htbp]
  \caption{TODO: Algorithm name}
  \label{alg:proposed}
  \begin{algorithmic}[1]
    \Require Initial state $\bx_0$, matrices $\bQ$, $\bR$, horizon $N$
    \Ensure Control sequence $\{\bu_k\}_{k=0}^{N-1}$
    \State Solve Riccati equation to obtain $\bP$
    \State Compute gain $\bK = (\bR + \bB\T\bP\bB)^{-1}\bB\T\bP\bA$
    \For{$k = 0, 1, \ldots, N-1$}
      \State $\bu_k \gets -\bK\,\bx_k$
      \State $\bx_{k+1} \gets \bA\,\bx_k + \bB\,\bu_k$
    \EndFor
    \State \Return $\{\bu_k\}$
  \end{algorithmic}
\end{algorithm}

\subsection{Stability Analysis}
% Theorem + Proof is standard for TAC/Automatica papers.
% For TIE/TIM, stability can be stated as a proposition if space is limited.

\begin{theorem}
  \label{thm:stability}
  Under Assumption~\ref{ass:stabilizable}, the closed-loop system
  \eqref{eq:system} with the proposed control law is asymptotically stable.
\end{theorem}
\begin{proof}
  TODO: Provide proof.
  Consider the Lyapunov function $V(\bx_k) = \bx_k\T\bP\,\bx_k > 0$.
  Then $\Delta V = V(\bx_{k+1}) - V(\bx_k) = \ldots < 0$. \qed
\end{proof}

\begin{remark}
  TODO: Add remark if needed.
\end{remark}


% =============================================================================
% SECTION V — SIMULATION
% =============================================================================
\section{Simulation Results}
\label{sec:simulation}
% MANDATORY: simulation/simulation_plan.md must exist before this section is written.
% All numbers in this section must match simulation/ CSV files exactly.

\subsection{Simulation Setup}
% Describe: software, hardware (CPU/GPU), simulation parameters, comparison baselines.

Simulations were conducted in MATLAB/Simulink (R202Xb) on a
\SI{3.2}{\giga\hertz} Intel Core i9 processor.
\cref{tab:sim_params} lists the simulation parameters.

\begin{table}[htbp]
  \caption{Simulation Parameters}
  \label{tab:sim_params}
  \centering
  % IEEEtran: caption ABOVE table, use \hline or booktabs
  \begin{tabular}{lll}
    \toprule
    Parameter & Symbol & Value \\
    \midrule
    Sampling period  & $T_s$ & \SI{10}{\milli\second} \\
    Simulation time  & $T$   & \SI{10}{\second} \\
    TODO: Add parameters & & \\
    \bottomrule
  \end{tabular}
\end{table}

\subsection{Comparison with Baselines}
% Compare with at least 2 baselines. Report quantitative metrics (RMSE, IAE, settling time, etc.)
% Numbers MUST match comparison/baseline_comparison.csv

\cref{tab:comparison} summarizes the performance comparison.
The proposed method achieves an RMSE of TODO~deg,
which is TODO\% lower than the PID baseline and TODO\% lower than Method~B \cite{TODO_ref}.

\begin{table}[htbp]
  \caption{Performance Comparison}
  \label{tab:comparison}
  \centering
  \begin{tabular}{lccc}
    \toprule
    Method & RMSE [deg] & Settling Time [s] & Overshoot [\%] \\
    \midrule
    PID (baseline)    & TODO & TODO & TODO \\
    Method B \cite{TODO_ref} & TODO & TODO & TODO \\
    \textbf{Proposed} & \textbf{TODO} & \textbf{TODO} & \textbf{TODO} \\
    \bottomrule
  \end{tabular}
\end{table}

\begin{figure}[htbp]
  \centering
  % \includegraphics[width=\columnwidth]{figures/sim_response.pdf}
  \rule{\columnwidth}{0.4pt}
  \vspace{1em}
  \caption{Simulation results: time-domain response comparison.
           (a) Proposed. (b) PID. (c) Method B.}
  \label{fig:sim_response}
\end{figure}


% =============================================================================
% SECTION VI — EXPERIMENTAL RESULTS
% =============================================================================
\section{Experimental Results}
\label{sec:experiment}
% MANDATORY: experiment/experimental_plan.md must exist before this section is written.
% All numbers in this section must match experiment/ data files exactly.

\subsection{Experimental Setup}
% Describe: hardware platform, sensors, actuators, data acquisition, safety limits.

TODO: Describe the experimental setup.

\begin{figure}[htbp]
  \centering
  % \includegraphics[width=\columnwidth]{figures/experimental_setup.jpg}
  \rule{\columnwidth}{0.4pt}
  \vspace{1em}
  \caption{Experimental setup. TODO: Add meaningful caption.}
  \label{fig:exp_setup}
\end{figure}

\subsection{Results}
% Present results in the same format as simulation for easy comparison.
% Report mean ± standard deviation over repeated trials.

TODO: Present experimental results.


% =============================================================================
% SECTION VII — DISCUSSION (optional, often merged into result sections)
% =============================================================================
% \section{Discussion}
% \label{sec:discussion}
% Discuss: Why does the proposed method work? Limitations? Future work hints?


% =============================================================================
% SECTION VIII — CONCLUSION
% =============================================================================
\section{Conclusion}
\label{sec:conclusion}
% Structure:
%   Paragraph 1: Restate the problem and the proposed approach (2–3 sentences)
%   Paragraph 2: Summarize key results with numbers
%   Paragraph 3: Limitations and future work

TODO: Write conclusion.

In this paper, we proposed TODO.
Simulation and experimental results demonstrated TODO\% improvement
in TODO metric compared to existing baselines.
Future work will investigate TODO.


% =============================================================================
% ACKNOWLEDGMENT
% =============================================================================
\section*{Acknowledgment}
% Standard funding acknowledgment format for Korean NRF/IITP grants:
% "This work was supported by [funding body] under Grant [number]."
This work was supported by the National Research Foundation of Korea (NRF)
under Grant NRF-TODO.
% K-UAM example:
% This work was supported by the Korea Research Institute for defense
% Technology planning and advancement (KRIT) under Grant KRIT-TODO
% (K-UAM Grand Challenge).


% =============================================================================
% APPENDIX (if needed)
% =============================================================================
% \appendices
% \section{Proof of Theorem~\ref{thm:stability}}
% \label{app:proof}
% TODO: Extended proof here.


% =============================================================================
% REFERENCES
% =============================================================================
% IEEEtran reference style:
%   - Numeric, sorted by order of appearance: [1], [2], ...
%   - Use \bibliographystyle{IEEEtran}
%   - Run: bibtex ieee_template

% Uncomment before last page to balance two-column layout:
% \balance

\bibliographystyle{IEEEtran}
\bibliography{references}

% =============================================================================
% BIOGRAPHY (required for some IEEE journals, optional for others)
% Skip for initial submission; add for final version if required.
% =============================================================================
% \begin{IEEEbiography}[{\includegraphics[width=1in,height=1.25in,clip,keepaspectratio]{photo_son.jpg}}]{Hungsun Son}
%   received the B.S., M.S., and Ph.D. degrees in mechanical engineering
%   from TODO University, in TODO, TODO, and TODO, respectively.
%   He is currently an Associate Professor with the Department of Mechanical
%   Engineering, Ulsan National Institute of Science and Technology (UNIST),
%   Ulsan, Republic of Korea.
%   His research interests include motion platforms, control systems,
%   and urban air mobility (UAM).
% \end{IEEEbiography}

\end{document}
% =============================================================================
% END OF TEMPLATE
% =============================================================================
